\documentclass[11pt,a4paper]{article}

% ── Packages ──────────────────────────────────────────────────────
\usepackage[utf8]{inputenc}
\usepackage[T1]{fontenc}
\usepackage{lmodern}
\usepackage[margin=2.5cm]{geometry}
\usepackage{graphicx}
\usepackage{xcolor}
\usepackage{listings}
\usepackage{booktabs}
\usepackage{longtable}
\usepackage{hyperref}
\usepackage{enumitem}
\usepackage{fancyhdr}
\usepackage{tcolorbox}
\usepackage{amsmath}

% ── Colour definitions ───────────────────────────────────────────
\definecolor{codegreen}{rgb}{0.13,0.55,0.13}
\definecolor{codegray}{rgb}{0.5,0.5,0.5}
\definecolor{codepurple}{rgb}{0.58,0,0.82}
\definecolor{backcolour}{rgb}{0.95,0.95,0.95}
\definecolor{warnred}{rgb}{0.8,0.1,0.1}

% ── Listings style ───────────────────────────────────────────────
\lstdefinestyle{code}{
    backgroundcolor=\color{backcolour},
    commentstyle=\color{codegreen},
    keywordstyle=\color{codepurple}\bfseries,
    stringstyle=\color{codegreen},
    basicstyle=\ttfamily\footnotesize,
    breakatwhitespace=false,
    breaklines=true,
    captionpos=b,
    keepspaces=true,
    showspaces=false,
    showstringspaces=false,
    showtabs=false,
    tabsize=2,
    frame=single,
    rulecolor=\color{codegray},
}
\lstset{style=code}

% ── tcolorbox styles ─────────────────────────────────────────────
\newtcolorbox{warningbox}{colback=red!5!white, colframe=red!75!black,
  title={\textbf{Warning}}, fonttitle=\bfseries}

\newtcolorbox{infobox}{colback=blue!5!white, colframe=blue!60!black,
  title={\textbf{Info}}, fonttitle=\bfseries}

\newtcolorbox{tipbox}{colback=green!5!white, colframe=green!60!black,
  title={\textbf{Tip}}, fonttitle=\bfseries}

% ── Header / Footer ─────────────────────────────────────────────
\pagestyle{fancy}
\fancyhf{}
\rhead{Collective Intelligence Project}
\lhead{LibSignal Team Instructions}
\cfoot{\thepage}

% ── Document ─────────────────────────────────────────────────────
\begin{document}

% ═══════════════════════════════════════════════════════════════════
\begin{titlepage}
\centering
\vspace*{3cm}
{\Huge\bfseries LibSignal Research\\[0.4cm] Team Instructions\par}
\vspace{1.5cm}
{\Large Collective Intelligence Project\par}
\vspace{1cm}
{\large Salman $\cdot$ Rashid $\cdot$ Madina\par}
\vspace{2cm}
{\large \today\par}
\vfill
{\normalsize
Repository: \url{https://github.com/sal0-h/LibSignalFork}\\[0.3cm]
Framework: LibSignal --- Reinforcement Learning for Traffic Signal Control}
\end{titlepage}

% ═══════════════════════════════════════════════════════════════════
\tableofcontents
\newpage

% ═══════════════════════════════════════════════════════════════════
\section{Overview}

LibSignal is an open-source reinforcement learning framework for \textbf{Traffic Signal Control (TSC)}.
It wraps traffic simulators (SUMO, CityFlow) behind a unified interface so that different RL agents can be trained and evaluated on various road networks without rewriting glue code.

We are using a \textbf{fork} of the original repository with small fixes for Python~3.11+ compatibility and realistic SUMO physics.

\subsection{Key Concepts}
\begin{description}[leftmargin=3cm, style=nextline]
    \item[Agent] The RL policy (e.g.\ DQN, PressLight, CoLight, MPLight).
    \item[World] The simulator backend (we use \textbf{SUMO}).
    \item[Network] The road topology and traffic demand (e.g.\ \texttt{sumo1x1}, \texttt{grid4x4}).
    \item[Episode] One full simulation run (default 3600\,s = 1 hour of simulated traffic).
    \item[Action Interval] How often the agent makes a decision (default every 10\,s).
\end{description}

% ═══════════════════════════════════════════════════════════════════
\section{Google Colab Setup}

All experiments run on Google Colab.  Each team member has a \textbf{personal output folder} on Google Drive so that checkpoints and logs are never lost when Colab disconnects.

\subsection{Step 1: Open the Notebook}

Open the shared Colab notebook.  The first cell is the \textbf{Setup Cell}.

\subsection{Step 2: Change Your Name}

At the top of the setup cell, change \texttt{MEMBER\_NAME} to your name:

\begin{lstlisting}[language=Python]
MEMBER_NAME = "salman"   # Options: "salman", "rashid", "madina"
\end{lstlisting}

This determines your personal save folder on Google Drive:
\begin{center}
\texttt{Collective Intelligence Project/colab\_\textit{<name>}/output\_data/}
\end{center}

\subsection{Step 3: Run the Setup Cell}

Click the play button.  The cell will:
\begin{enumerate}
    \item Mount Google Drive.
    \item Create your personal output folder (if it doesn't exist).
    \item Install SUMO (first run only, takes $\sim$2 minutes).
    \item Clone or update the LibSignal fork.
    \item Install Python dependencies.
    \item Create a \textbf{symlink} from the local \texttt{data/output\_data} directory to your Drive folder, so all logs and model checkpoints are automatically saved to the cloud.
\end{enumerate}

\begin{warningbox}
You \textbf{must} run the setup cell every time you open a new Colab session.
Colab resets the virtual machine when the session ends, so the symlink and SUMO installation are lost.
\end{warningbox}

\subsection{Where Do Outputs Go?}

When you run an experiment, outputs are saved to:
\begin{center}
\texttt{data/output\_data/tsc/\{world\}\_\{agent\}/\{network\}/\{prefix\}/}
\end{center}

Because of the symlink, this maps directly to your Google Drive:
\begin{center}
\texttt{Drive/.../colab\_\textit{<name>}/output\_data/tsc/sumo\_presslight/grid4x4/test/}
\end{center}

Inside that folder you will find:
\begin{itemize}
    \item \texttt{logger/} --- training logs (one \texttt{.log} file per run).
    \item \texttt{model/} --- saved model checkpoints (\texttt{.pt} files, saved every 5 episodes by default).
\end{itemize}

% ═══════════════════════════════════════════════════════════════════
\section{Running Experiments}

There are two ways to run experiments.

% ── Option 1 ─────────────────────────────────────────────────────
\subsection{Option 1: Edit the Config File Directly (Recommended)}

This is the most reliable method.

\begin{enumerate}
    \item Open the Colab \textbf{file explorer} (folder icon on the left sidebar).
    \item Navigate to \texttt{LibSignalFork/configs/tsc/}.
    \item Open the YAML file for your agent (e.g.\ \texttt{presslight.yml}).
    \item Edit the hyperparameters to match the values from the paper.
    \item Save the file (\texttt{Ctrl+S}).
    \item Run the experiment cell:
\end{enumerate}

\begin{lstlisting}[language=bash]
!cd /content/LibSignalFork && python run.py \
    --task tsc \
    --agent presslight \
    --world sumo \
    --network grid4x4 \
    --prefix my_experiment_name \
    --dataset onfly
\end{lstlisting}

% ── Option 2 ─────────────────────────────────────────────────────
\subsection{Option 2: Smart Config Injector Cell}

The notebook includes a ``Smart Config Injector'' cell that programmatically modifies the YAML config before each run and restores it afterwards.  Edit the \texttt{OVERRIDES} dictionary in that cell.

\begin{lstlisting}[language=Python]
OVERRIDES = {
    "model": {
        "learning_rate": 0.001,
        "batch_size": 64,
        "epsilon": 0.1,
        "epsilon_decay": 0.995,
    },
    "trainer": {
        "episodes": 200,
        "test_steps": 3600,
        "buffer_size": 10000,
    }
}
\end{lstlisting}

\begin{tipbox}
Use the \texttt{PREFIX} variable to give each run a unique name.
This prevents experiments from overwriting each other.
Example: \texttt{PREFIX = "lr001\_eps01\_bs64"}.
\end{tipbox}

% ═══════════════════════════════════════════════════════════════════
\section{Command-Line Arguments Reference}

All arguments are passed to \texttt{run.py}:

\begin{center}
\begin{longtable}{llp{7cm}}
\toprule
\textbf{Flag} & \textbf{Default} & \textbf{Description} \\
\midrule
\endhead
\texttt{-{}-task} / \texttt{-t}       & \texttt{tsc}       & Task type. Always \texttt{tsc} for us. \\
\texttt{-{}-agent} / \texttt{-a}      & \texttt{dqn}       & Agent name (must match a YAML file in \texttt{configs/tsc/}). \\
\texttt{-{}-world} / \texttt{-w}      & \texttt{cityflow}  & Simulator. Use \texttt{sumo}. \\
\texttt{-{}-network} / \texttt{-n}    & \texttt{cityflow1x1} & Network config name (see Section~\ref{sec:networks}). \\
\texttt{-{}-dataset} / \texttt{-d}    & \texttt{onfly}     & Dataset type. Always \texttt{onfly}. \\
\texttt{-{}-prefix}                    & \texttt{test}      & Output subfolder name. \textbf{Use this to tag experiments.} \\
\texttt{-{}-seed}                      & \texttt{None}      & Random seed for reproducibility. \\
\texttt{-{}-ngpu}                      & \texttt{-1}        & GPU id (\texttt{-1} = CPU, \texttt{0} = first GPU). \\
\texttt{-{}-interface}                 & \texttt{libsumo}   & \texttt{libsumo} (fast) or \texttt{traci} (slow, supports GUI). \\
\texttt{-{}-delay\_type}              & \texttt{apx}       & \texttt{apx} (approximate) or \texttt{real} delay calculation. \\
\bottomrule
\end{longtable}
\end{center}

\begin{warningbox}
Hyperparameters like \texttt{learning\_rate}, \texttt{epsilon}, \texttt{batch\_size}, \texttt{episodes}, etc.\ \textbf{cannot} be set from the command line.
They must be set in the YAML config files or injected programmatically.
\end{warningbox}

% ═══════════════════════════════════════════════════════════════════
\section{Available Agents}

\begin{center}
\begin{longtable}{lll}
\toprule
\textbf{Agent Name} & \textbf{Config File} & \textbf{Type} \\
\midrule
\endhead
\texttt{maxpressure}  & \texttt{maxpressure.yml}  & Heuristic baseline (no learning) \\
\texttt{fixedtime}    & \texttt{fixedtime.yml}    & Fixed-time baseline (no learning) \\
\texttt{sotl}         & \texttt{sotl.yml}         & Self-Organizing Traffic Light \\
\texttt{dqn}          & \texttt{dqn.yml}          & Deep Q-Network \\
\texttt{presslight}   & \texttt{presslight.yml}   & PressLight (pressure-based DQN) \\
\texttt{mplight}      & \texttt{mplight.yml}      & MPLight (FRAP-based, pressure) \\
\texttt{colight}      & \texttt{colight.yml}      & CoLight (graph attention network) \\
\texttt{ppo}          & \texttt{ppo.yml}          & Proximal Policy Optimization \\
\texttt{frap}         & \texttt{frap.yml}         & FRAP \\
\texttt{maddpg}       & \texttt{maddpg.yml}       & Multi-Agent DDPG \\
\bottomrule
\end{longtable}
\end{center}

% ═══════════════════════════════════════════════════════════════════
\section{Available SUMO Networks}
\label{sec:networks}

These are the \texttt{-{}-network} values you can use with \texttt{-{}-world sumo}:

\begin{center}
\begin{longtable}{lp{8cm}}
\toprule
\textbf{Network Name} & \textbf{Description} \\
\midrule
\endhead
\texttt{sumo1x1}          & Synthetic single intersection \\
\texttt{sumo1x3}          & Synthetic 1$\times$3 arterial \\
\texttt{sumo4x4}          & Synthetic 4$\times$4 grid \\
\texttt{sumo1x21}         & Synthetic 1$\times$21 arterial \\
\texttt{sumo7x28}         & Large 7$\times$28 grid \\
\texttt{sumohz1x1}        & Real-world Hangzhou single intersection \\
\texttt{sumohz4x4}        & Real-world Hangzhou 4$\times$4 grid \\
\texttt{grid4x4}          & Alias for synthetic 4$\times$4 grid (CityFlow-style name) \\
\texttt{cologne1}         & Real-world Cologne single intersection \\
\texttt{cologne3}         & Real-world Cologne 1$\times$3 corridor \\
\bottomrule
\end{longtable}
\end{center}

\begin{infobox}
For multi-intersection agents like \textbf{CoLight} and \textbf{MPLight}, use multi-intersection networks
(\texttt{grid4x4}, \texttt{sumo4x4}, \texttt{cologne3}, etc.).
Single-intersection networks work best for \texttt{dqn}, \texttt{presslight}, \texttt{ppo}, etc.
\end{infobox}

% ═══════════════════════════════════════════════════════════════════
\section{Configuration (YAML) Reference}

Each agent has a YAML file in \texttt{configs/tsc/} that \texttt{includes} the base config.
Agent-specific parameters override the base defaults.
The inheritance chain is:

\begin{center}
\texttt{configs/tsc/base.yml} $\longrightarrow$ \texttt{configs/tsc/<agent>.yml}
\end{center}

\subsection{Key Trainer Parameters}

\begin{center}
\begin{longtable}{llp{7cm}}
\toprule
\textbf{Parameter} & \textbf{Default} & \textbf{Description} \\
\midrule
\endhead
\texttt{episodes}          & 200   & Number of training episodes. \\
\texttt{steps}             & 3600  & Simulation steps per training episode (seconds). \\
\texttt{test\_steps}       & 3600  & Simulation steps per test episode. \\
\texttt{action\_interval}  & 10    & Agent makes a decision every $N$ seconds. \\
\texttt{yellow\_length}    & 5     & Yellow phase duration (seconds). \\
\texttt{learning\_start}   & 5000  & Steps before learning begins (replay buffer warm-up). \\
\texttt{buffer\_size}      & 5000  & Replay buffer capacity. \\
\texttt{update\_model\_rate}  & 1  & Update model every $N$ episodes. \\
\texttt{update\_target\_rate} & 10 & Update target network every $N$ episodes. \\
\texttt{test\_when\_train}  & True & Run a test episode after each training episode. \\
\bottomrule
\end{longtable}
\end{center}

\subsection{Key Model (Agent) Parameters}

\begin{center}
\begin{longtable}{llp{7cm}}
\toprule
\textbf{Parameter} & \textbf{Default} & \textbf{Description} \\
\midrule
\endhead
\texttt{learning\_rate}   & 0.001 & Optimizer learning rate. \\
\texttt{batch\_size}      & 64    & Mini-batch size for training. \\
\texttt{gamma}            & 0.95  & Discount factor $\gamma$. \\
\texttt{epsilon}          & 0.5   & Initial exploration rate $\varepsilon$. \\
\texttt{epsilon\_decay}   & 0.99  & Multiplicative decay per episode. \\
\texttt{epsilon\_min}     & 0.05  & Minimum $\varepsilon$. \\
\texttt{grad\_clip}       & 5.0   & Gradient clipping threshold. \\
\texttt{train\_model}     & False & Set to \texttt{True} to enable training. \\
\texttt{load\_model}      & False & Load a previously saved checkpoint. \\
\bottomrule
\end{longtable}
\end{center}

\subsection{Logger Parameters}

\begin{center}
\begin{longtable}{llp{7cm}}
\toprule
\textbf{Parameter} & \textbf{Default} & \textbf{Description} \\
\midrule
\endhead
\texttt{save\_rate}  & 5    & Save model checkpoint every $N$ episodes. \\
\texttt{save\_model} & True & Whether to save checkpoints at all. \\
\bottomrule
\end{longtable}
\end{center}

% ═══════════════════════════════════════════════════════════════════
\section{Team Member Assignments}

\begin{center}
\begin{tabular}{llll}
\toprule
\textbf{Member} & \textbf{Agent(s)} & \textbf{Network} & \textbf{Drive Folder} \\
\midrule
Salman  & PressLight (baselines) & grid4x4  & \texttt{colab\_salman/} \\
Rashid  & CoLight (graph models) & grid4x4  & \texttt{colab\_rashid/} \\
Madina  & MPLight (comparison)   & cologne3 & \texttt{colab\_madina/} \\
\bottomrule
\end{tabular}
\end{center}

Each member should use the \texttt{PREFIX} flag to separate different runs:
\begin{lstlisting}[language=bash]
--prefix run1_lr001_eps01
--prefix run2_lr0001_eps08
--prefix final_200ep
\end{lstlisting}

% ═══════════════════════════════════════════════════════════════════
\section{Verifying a Run}

\begin{enumerate}
    \item Run the setup cell.
    \item Start a training run.
    \item Wait $\sim$2 minutes for the first episode to finish.
    \item Open Google Drive in a new tab.
    \item Navigate to: \texttt{Collective Intelligence Project} $\to$ \texttt{colab\_\textit{<name>}} $\to$ \texttt{output\_data}.
    \item You should see a folder tree like:
\begin{lstlisting}[language=bash]
output_data/
  tsc/
    sumo_presslight/
      grid4x4/
        test/             # <-- your PREFIX
          logger/
            2026_02_10-16_39_57_BRF.log
          model/
            5_0.pt        # checkpoint at episode 5
            10_0.pt       # checkpoint at episode 10
            ...
\end{lstlisting}
    \item If you see the \texttt{model/} folder appearing, your setup is working correctly.
\end{enumerate}

% ═══════════════════════════════════════════════════════════════════
\section{Important Warnings}

\begin{warningbox}
\textbf{Colab will disconnect if you close the tab.}
There is no way around this.  Keep the tab open and prevent your computer from sleeping.
\end{warningbox}

\subsection{Anti-Disconnect Hack}

To prevent Colab from disconnecting due to idle timeout (even while training is running), open the browser console (\texttt{Ctrl+Shift+I} $\to$ Console tab) and paste:

\begin{lstlisting}[language=Java]
function ClickConnect(){
    console.log("Keeping Colab alive...");
    document.querySelector("colab-toolbar-button#connect").click()
}
setInterval(ClickConnect, 60000)
\end{lstlisting}

This clicks the ``Connect'' button every 60 seconds to prevent idle disconnection.

\subsection{Other Gotchas}

\begin{itemize}
    \item \textbf{Do not} run \texttt{git pull} in the Colab repo directory if you have manually edited config files---it will overwrite your changes. Use the Smart Config Injector instead, which restores the original config after each run.
    \item \textbf{Do not} change \texttt{world.dir} in \texttt{base.yml}. It must remain \texttt{data/} for the output symlink to work.
    \item Colab free tier has a \textbf{$\sim$12 hour} session limit. For 200 episodes on a 4$\times$4 grid, this should be enough, but monitor your progress.
    \item If you see \texttt{SUMO\_HOME} errors, the setup cell didn't run. Re-run it.
\end{itemize}

% ═══════════════════════════════════════════════════════════════════
\section{Output Metrics}

The training log prints the following metrics each episode:

\begin{center}
\begin{longtable}{lp{9cm}}
\toprule
\textbf{Metric} & \textbf{Meaning} \\
\midrule
\endhead
\texttt{q\_loss}        & The Q-network loss (should decrease over training). \\
\texttt{rewards}        & Average reward per step (higher is better). \\
\texttt{queue}          & Average queue length at intersections (lower is better). \\
\texttt{delay}          & Average vehicle delay (lower is better). \\
\texttt{throughput}     & Number of vehicles that completed their trip (higher is better). \\
\texttt{real avg travel time} & Mean travel time of all vehicles (lower is better). This is the \textbf{primary metric} for comparison. \\
\bottomrule
\end{longtable}
\end{center}

% ═══════════════════════════════════════════════════════════════════
\section{Quick Reference Card}

\begin{tcolorbox}[colback=gray!5!white, colframe=gray!75!black, title=\textbf{Cheat Sheet}]
\begin{lstlisting}[language=bash]
# Baseline (no learning, instant)
python run.py -a maxpressure -w sumo -n sumo1x1 --prefix baseline

# DQN on single intersection
python run.py -a dqn -w sumo -n sumo1x1 --prefix dqn_run1

# PressLight on 4x4 grid
python run.py -a presslight -w sumo -n grid4x4 --prefix pl_run1

# CoLight on 4x4 grid
python run.py -a colight -w sumo -n grid4x4 --prefix col_run1

# MPLight on Cologne corridor
python run.py -a mplight -w sumo -n cologne3 --prefix mpl_run1
\end{lstlisting}
\textbf{All commands assume}: \texttt{-{}-task tsc -{}-dataset onfly}
\end{tcolorbox}

\end{document}
